\documentclass[11pt]{article}
\usepackage[T1]{fontenc}
\usepackage{lmodern}
\usepackage{amsmath,amssymb,amsthm}
\usepackage[margin=1in]{geometry}
\usepackage{microtype}
\usepackage{hyperref}
\hypersetup{colorlinks=true,linkcolor=blue,citecolor=blue,urlcolor=blue,
  pdftitle={Recurrent constant-term words for a degenerate Laurent polynomial},
  pdfauthor={Xiaojian Zeng}}

\newtheorem{theorem}{Theorem}
\newtheorem{lemma}[theorem]{Lemma}
\newtheorem{proposition}[theorem]{Proposition}
\theoremstyle{remark}
\newtheorem{remark}[theorem]{Remark}
\DeclareMathOperator{\Res}{Res}
\DeclareMathOperator{\FS}{FS}
\DeclareMathOperator{\supp}{supp}
\newcommand{\Nzero}{\mathbb N_0}
\newcommand{\Fp}{\mathbb F_p}
\title{Recurrent constant-term words for\\$x^{p-1}+x^{-1}$}
\author{Xiaojian Zeng\\
  \small PBC School of Finance, Tsinghua University, 100190\\
  \small\texttt{xiaojianzeng@foxmail.com}}
\date{}

\begin{document}
\maketitle
\begin{abstract}
Let $p$ be any prime, put $P(x)=x^{p-1}+x^{-1}$, and let
$Q\in\Fp[x,x^{-1}]$. We prove that the word
$u_n=[x^0]P(x)^nQ(x)$, indexed from $n=0$, is recurrent if and only if it
is identically zero. Here recurrence means that every occurring finite
contiguous block occurs arbitrarily far out, without a bound on return gaps.
We also determine all such $Q$: they are exactly $P F(P)$ with
$F\in\Fp[t]$. The proof identifies the generating functions of these words
with polynomials in a lacunary series and computes the kernel of the
constant-term map. A digit-sum obstruction then excludes the translated
finite-sums sets forced by recurrence. This settles the specified example
in Kohen's question on completely degenerate constant-term sequences.
\end{abstract}

\section{Statement and context}

Write $[x^0]$ for the constant coefficient of a Laurent polynomial.
Throughout, $p$ is prime, $K=\Fp$, and
\begin{equation}\label{eq:family}
 P(x)=x^{p-1}+x^{-1},\qquad
 u_n(Q)=[x^0]\bigl(P(x)^nQ(x)\bigr)\quad(n\in\Nzero).
\end{equation}
A word $v=(v_n)_{n\geq0}$ is \emph{recurrent} if, for every
$i,N\geq0$ and $\ell\geq1$, there is $j>N$ such that
$v_{j+k}=v_{i+k}$ for $0\leq k<\ell$. This agrees with
\cite[Definition~18]{Kohen2026}. In particular, the initial term
$u_0(Q)=[x^0]Q$ is included.

\begin{theorem}\label{thm:main}
For every prime $p$ and every $Q\in K[x,x^{-1}]$, the following are equivalent:
\begin{enumerate}
\item[(i)] The word $u(Q)$ defined in \eqref{eq:family} is recurrent.
\item[(ii)] $u_n(Q)=0$ for every $n\geq0$.
\item[(iii)] $Q=P F(P)$ for some polynomial $F\in K[t]$.
\end{enumerate}
\end{theorem}

The space in (iii) is $P K[P]$, the polynomials in $P$ with zero constant
coefficient as polynomials in $P$. It is a $K$-linear subspace, and should
not be replaced by the Laurent-ring ideal $P K[x,x^{-1}]$.
The statement includes $p=2$, $Q=0$, and arbitrary finite positive and
negative Laurent exponents.

Constant-term sequences modulo prime powers admit finite-automaton
representations by the construction of Rowland and Zeilberger
\cite{RowlandZeilberger2014}. Kohen's earlier work treats uniform recurrence
for symmetric Laurent trinomials; the relevant general-$Q$ statement
assumes $p>2$ \cite[Proposition~15]{Kohen2025}.
His subsequent general result shows that if an unweighted constant term
vanishes modulo $p$, then zeros have density one in every weighted sequence,
and a nonzero such word is not uniformly recurrent
\cite[Theorem~14]{Kohen2026}. Ordinary recurrence still holds if some
positive unweighted term is nonzero \cite[Proposition~19]{Kohen2026}.
Question~20 of that paper asks about the remaining case in which every
positive unweighted term vanishes, explicitly giving the polynomial
\eqref{eq:family} as an example.

Theorem~\ref{thm:main} classifies ordinary recurrence and the entire zero
kernel for this example. The inspected results above do not supply this
classification. This comparison is limited to the primary material checked
through September~12, 2026; it is not an assertion of exhaustive priority.
The argument makes no claim about other degenerate $P$ or moduli $p^a$ with
$a>1$, so the generality of Kohen's question is not resolved here.

The proof uses a constraint stronger than density-zero support. We first
reduce every generating function to a polynomial in one lacunary series.
We then show that its support cannot contain the translated finite-sums set
that any recurrent nonzero word must contain. These elementary arguments
are given in full; no separate originality is asserted for their components.

\section{The generating-function map and its kernel}

Set
\begin{equation}\label{eq:Y}
 a_j=\frac{p^j-1}{p-1}\quad(j\geq1),\qquad
 Y(z)=\sum_{j\geq1}z^{a_j},\qquad
 \Phi(Q)=\sum_{n\geq0}u_n(Q)z^n.
\end{equation}
The integer identity $a_{j+1}=p a_j+1$ and Frobenius give
\begin{equation}\label{eq:inverse}
 Y=z(1+Y^p),\qquad Y=z+O(z^2).
\end{equation}
Thus $Y$ is the compositional inverse of $f(x)=x/(1+x^p)$.

\begin{lemma}\label{lem:monomial}
For every integer $r\geq0$, one has $\Phi(x^r)=Y^r$.
\end{lemma}
\begin{proof}
We use formal residues: $\Res_x G(x)\,dx=[x^{-1}]G(x)$.
The change-of-variable identity is
\begin{equation}\label{eq:residue}
 \Res_z G(z)\,dz=\Res_x G(f(x))f'(x)\,dx
 \qquad\bigl(G\in K((z)),\ f=x+O(x^2)\bigr).
\end{equation}
For its characteristic-zero form, see
\cite[Theorem~4.1.1]{Gessel2016}. To justify its use over $K$, lift the
coefficients of $G$ and $f$ to integers, keeping the linear coefficient of
$f$ equal to $1$. Only finitely many coefficients affect either residue.
Over $\mathbb Q$, linearity reduces the identity to
$G(z)=z^m$: for $m\ne-1$, the residue of
$f^m f'=(f^{m+1})'/(m+1)$ is zero; for $m=-1$, writing $f=xh$ with
$h\in1+x\mathbb Z[[x]]$ gives $f'/f=x^{-1}+h'/h$, of residue $1$.
Both residues are integral expressions, so reducing this identity modulo
$p$ proves \eqref{eq:residue}. No division by an index takes place in $K$.

In $K[[x]]$ we have $f'(x)=(1+x^p)^{-1}$ and $Y(f(x))=x$.
For $r,n\geq0$, substitution gives
\begin{align*}
 [z^n]Y(z)^r
 &=\Res_x x^r\left(\frac{x}{1+x^p}\right)^{-n-1}
                  \frac{dx}{1+x^p}\\
 &=[x^{n-r}](1+x^p)^n
  =[x^0]\bigl(P(x)^n x^r\bigr).
\end{align*}
Equality of every coefficient proves the lemma.
\end{proof}

In particular, $\Phi(1)=1$, so $[x^0]P^n=0$ for all $n>0$.
The next proposition includes all negative exponents and all cancellations
among Laurent monomials.

\begin{proposition}\label{prop:kernel}
The linear map $\Phi:K[x,x^{-1}]\longrightarrow K[[z]]$ has image $K[Y]$
and kernel $P K[P]$.
\end{proposition}
\begin{proof}
Consider the quotient
\[
 R=K[t,X]/(X^p-tX+1).
\]
In $R$, one has $X(t-X^{p-1})=1$. Hence the homomorphisms
\begin{align*}
 R&\longrightarrow K[x,x^{-1}], & X&\longmapsto x,&t&\longmapsto P,\\
 K[x,x^{-1}]&\longrightarrow R, & x&\longmapsto X,&x^{-1}&\longmapsto t-X^{p-1}
\end{align*}
are well defined and inverse: their compositions fix the displayed
generators, with $X^{p-1}+X^{-1}=t$ in $R$.
Division by the monic polynomial in $X$ makes $R$ a free $K[t]$-module
with basis $1,X,\ldots,X^{p-1}$. Consequently every Laurent polynomial
has a unique expression
\begin{equation}\label{eq:module}
 Q=\sum_{r=0}^{p-1}f_r(P)x^r,\qquad f_r\in K[t].
\end{equation}

The map $\Phi$ is linear, but no multiplicativity is used. Instead, shifting
the index in its definition gives
\begin{equation}\label{eq:shift}
 \Phi(PQ)=\frac{\Phi(Q)-[x^0]Q}{z}.
\end{equation}
If $\Phi(Q)=H(Y)$ for $H\in K[y]$, the constant coefficient is
$[x^0]Q=H(0)$. Since \eqref{eq:inverse} implies
$z^{-1}=Y^{-1}+Y^{p-1}$, define
\begin{equation}\label{eq:L}
 L(H)(y)=(H(y)-H(0))(y^{-1}+y^{p-1});
 \qquad \Phi(PQ)=L(H)(Y).
\end{equation}
The expression $L(H)$ is polynomial because $y$ divides $H(y)-H(0)$.
The operator kills constants. If $H$ has positive degree $d$, then $L(H)$
has degree $d+p-1$ and the same leading coefficient.
Lemma~\ref{lem:monomial}, \eqref{eq:module}, and \eqref{eq:L} now imply
$\Phi(Q)\in K[Y]$ for every Laurent $Q$.

The $K$-basis $\{P^k x^r:k\geq0,\ 0\leq r\leq p-1\}$ of the Laurent
ring has zero images for $r=0$, $k\geq1$. Its remaining polynomial images are
\[
 1,\qquad L^k(y^r)\quad(k\geq0,\ 1\leq r\leq p-1).
\]
These are monic of pairwise distinct degrees $0$ and $r+k(p-1)$,
which exhaust the nonnegative integers. They are therefore a basis of
$K[y]$, by elimination of leading terms.
Substitution $H\mapsto H(Y)$ is injective: the lowest nonzero monomial of
$H$ has the same degree and coefficient after $Y=z+O(z^2)$ is substituted.
It follows that $\Phi(Q)=0$ precisely when $f_r=0$ for $r>0$ and
$f_0(0)=0$ in \eqref{eq:module}. This is exactly $Q\in P K[P]$,
and the same basis argument proves surjectivity onto $K[Y]$.
\end{proof}

\begin{remark}\label{rem:characteristic}
The full relation in the quotient above matters. Over $K(t)$ its defining
polynomial has derivative $-t$. This is the family $X^p+cX+d$ for which
Zeng's preprint \cite{Zeng_2026} records in Proposition~6.1 that the normalized
Bézout pair for the polynomial and its derivative is $(0,c^{-1})$,
independently of $d$. Thus these derivative data cannot recover the
constant coefficient. Our explicit presentation retains that coefficient:
the value $1$ is what gives the inverse $X^{-1}=t-X^{p-1}$.
\end{remark}

\section{Recurrence and a digit-sum obstruction}

For positive integers $t_1,t_2,\ldots$, write
\[
 \FS_0(t_i)=\left\{\sum_{i\in I}t_i:
          I\subset\mathbb N\text{ finite}\right\},
\]
including the empty sum. Finite-sums structure also occurs in the study of
sparse automatic sets. Byszewski and Konieczny
\cite[Definition~3.11 and Theorem~3.13]{ByszewskiKonieczny2020}
formulate an arid-or-IPS alternative, where the IPS notion allows shifts
depending on the stage. Their discussion following that theorem
distinguishes this from a single translate of an IP set. Here we need one
fixed translate, and establish both necessary lemmas directly.

\begin{lemma}\label{lem:returns}
If a word $v$ is recurrent and $v_a=c$, there exist positive integers $t_i$
such that $v_{a+s}=c$ for every $s\in\FS_0(t_i)$.
\end{lemma}
\begin{proof}
Suppose $t_1,\ldots,t_m$ have been chosen so that the assertion holds for
their subset sums. Put $T_m=\sum_{i=1}^m t_i$, with $T_0=0$.
Recurrence gives an occurrence of the prefix $v_0\cdots v_{a+T_m}$
starting at an index $t_{m+1}>T_m$. For each old subset sum $s$, this
copies $v_{a+s}=c$ to $v_{a+s+t_{m+1}}=c$ while retaining all old
positions. Induction proves the assertion, since every finite subset
is included at a finite stage. No bound on return gaps is required.
\end{proof}

Let $s_p(m)$ denote the sum of the base-$p$ digits of $m\in\Nzero$.
For an integer $D\geq0$, define
\begin{equation}\label{eq:SD}
 S_{p,D}=\{n\geq0:s_p((p-1)n+r)\leq D
                   \text{ for some }r\in\{0,\ldots,D\}\}.
\end{equation}

\begin{lemma}\label{lem:digits}
For every $a\geq0$ and every infinite sequence of positive integers $t_i$,
one has $a+\FS_0(t_i)\not\subseteq S_{p,D}$.
\end{lemma}
\begin{proof}
For any $K\geq0$, every unused tail of $(t_i)$ contains a nonempty
consecutive block whose sum is divisible by $p^K$. Indeed, two among its
first $p^K+1$ partial sums, including the empty sum, agree modulo $p^K$.
Their difference is the sum of such a block and is positive.

Choose $K_1\geq0$ with $p^{K_1}>(p-1)a+D$, and take a block sum $b_1$
divisible by $p^{K_1}$. After choosing $b_i$, choose $K_{i+1}$ with
$p^{K_{i+1}}>(p-1)b_i$, and take a block sum $b_{i+1}$ divisible by
$p^{K_{i+1}}$ from a tail strictly after the indices already used.
Continue until $b_1,\ldots,b_{D+1}$ have been chosen. The blocks are
disjoint, so their total belongs to $\FS_0(t_i)$.

Since $b_i\geq p^{K_i}$, we have $K_{i+1}>K_i$.
The digits of $(p-1)b_i$ start at position $K_i$ or later and, for
$i\leq D$, end below position $K_{i+1}$. For every $0\leq r\leq D$,
the digits of $(p-1)a+r$ lie below $K_1$.
These summands therefore occupy disjoint digit positions. For the single
index $n=a+b_1+\cdots+b_{D+1}$ there are no carries between them, and
simultaneously for all $0\leq r\leq D$,
\begin{align*}
 s_p((p-1)n+r)
 &=s_p((p-1)a+r)+\sum_{i=1}^{D+1}s_p((p-1)b_i)\\
 &\geq D+1.
\end{align*}
Every positive $b_i$ contributes at least one to the sum.
Thus $n\in a+\FS_0(t_i)$ and $n\notin S_{p,D}$, as required.
This argument includes $D=0$ and $p=2$.
\end{proof}

\begin{lemma}\label{lem:support}
If $0\ne H\in K[y]$ has degree $D$, then
$\supp H(Y)\subseteq S_{p,D}$, where the support consists of indices
with nonzero coefficient.
\end{lemma}
\begin{proof}
If $[z^n]H(Y)\ne0$, at least one $Y^r$, $0\leq r\leq D$, contributes
a nonzero coefficient there. For $r>0$ there must be an ordered tuple
$j_1,\ldots,j_r\geq1$ with $n=a_{j_1}+\cdots+a_{j_r}$; only finitely
many tuples contribute at any fixed index. Hence
\[
 (p-1)n+r=p^{j_1}+\cdots+p^{j_r},\qquad
 s_p((p-1)n+r)\leq r\leq D.
\]
Carrying in base $p$ decreases digit sum, which proves the inequality.
For $r=0$ only $n=0$ occurs and satisfies the same condition.
Cancellations among coefficients can only remove support points.
\end{proof}

\begin{proof}[Proof of Theorem~\ref{thm:main}]
By Proposition~\ref{prop:kernel}, write $\Phi(Q)=H(Y)$.
If $H\ne0$ and $D=\deg H$, injectivity of substitution gives a nonzero
coefficient $u_a(Q)=c$.
If $u(Q)$ were recurrent, Lemma~\ref{lem:returns} would put some
$a+\FS_0(t_i)$ in the support of $H(Y)$.
Lemma~\ref{lem:support} puts this support in $S_{p,D}$, contradicting
Lemma~\ref{lem:digits}. Thus recurrence forces $H=0$, or equivalently
$u(Q)=0$. The zero word is recurrent, and Proposition~\ref{prop:kernel}
identifies its preimage as $P K[P]$.
\end{proof}

\section{Boundary cases and reproducibility}

For $Q=1$ the word is $(1,0,0,\ldots)$ and fails recurrence, whereas
$Q=P$ gives the zero word, including its initial term.
For $Q=x$, the support is exactly $\{a_j:j\geq1\}$; its successive gaps
are $p,p^2,p^3,\ldots$. These examples illustrate why the indexing from
zero and the precise algebraic kernel are necessary.

Sparsity by itself would not prove the theorem. For example, let $W_0=1$
and $W_{m+1}=W_m0^{4|W_m|}W_m$. The nested prefixes define a recurrent
infinite word: every finite block lies in some $W_m$ and is copied
arbitrarily far out in later second copies. Since $|W_m|=6^m$ and
$W_m$ contains $2^m$ ones, the number of ones in a prefix of length
$6^m\leq N<6^{m+1}$ is at most $2^{m+1}$, giving density zero.
The simultaneous digit constraint \eqref{eq:SD} supplies the additional
obstruction used above.

The accompanying supplementary programs use exact Python integer
arithmetic and finite-field reduction. The binomial identity
\begin{equation}\label{eq:binomial}
 [x^0]P^n x^e=
 \begin{cases}
 \displaystyle\binom{n}{(n-e)/p}\bmod p,
    & (n-e)/p\in\{0,\ldots,n\},\\
 0,&\text{otherwise},
 \end{cases}
\end{equation}
follows by expanding $P^n=x^{-n}(1+x^p)^n$.
One implementation evaluates these binomial coefficients by Lucas's digit
congruence \cite[Theorem~1.1]{LaugierSaikia2012} and computes images of
negative monomials from
$P x^{1-m}=x^{p-m}+x^{-m}$. It checks all $236{,}775$ coefficient
comparisons for $p\in\{2,3,5,7,11,13,17\}$, $-16\leq e\leq16$,
and $0\leq n\leq1024$, with $29{,}799$ direct integer-binomial
crosschecks through $n=128$.

A second implementation uses literal repeated Laurent multiplication
and the quotient normal form \eqref{eq:module}. It checks all $502{,}250$
coefficient comparisons for
$p\in\{2,3,5,7,11,13,17,19,23,29\}$, $-24\leq e\leq24$,
and $0\leq n\leq1024$. For each of these ten primes and
$R\in\{1,2,4,8,16,24\}$, it also checks the kernel of the matrix
$([x^0]P^n x^e)_{0\leq n\leq R(p-1),\,-R\leq e\leq R}$;
the computed nullity is $\lfloor R/(p-1)\rfloor$, with basis given by the
coefficient vectors of $P^k$, $1\leq k\leq\lfloor R/(p-1)\rfloor$.
The two programs respectively construct and check $207$ finite instances
of the digit-block argument, including $D=0$ and degrees through $16$.
The precise generator families and ranges are specified in the supplement.
All stated checks pass. They test finite identities and matrices; the
infinite recurrence classification follows from the written proof.

\paragraph{Computational and AI assistance.}
Language-model assistance was used for literature screening, proof
development, counterexample checks, program construction, a separate
mathematical review pass, and manuscript preparation. The second program
was developed in that separate pass without importing the first program.
This describes distinct implementations and review contexts, not statistical
independence or external peer review. No proof assistant was used.
The supplementary programs are diagnostic and are not premises of
Theorem~\ref{thm:main}.

\bibliographystyle{plainurl}
\bibliography{references}
\end{document}
